\documentclass[11pt]{article}
\pdfobjcompresslevel=0 % force classic uncompressed xref table (avoids MiKTeX broken-xref-stream bug)
\usepackage[margin=2.3cm]{geometry}
\usepackage[T1]{fontenc}\usepackage{lmodern}
\usepackage{amsmath,amssymb,mathtools}
\usepackage{tcolorbox}\tcbuselibrary{skins,breakable}
\usepackage{booktabs}\usepackage{array}\usepackage{enumitem}
\newcolumntype{L}[1]{>{\raggedright\arraybackslash}p{#1}}
\usepackage{fancyhdr}\setlength{\headheight}{14.5pt}
\usepackage{parskip}\usepackage{xcolor}\usepackage{hyperref}
\usepackage{tikz}\usetikzlibrary{arrows.meta,positioning}

\definecolor{dtfblue}{HTML}{1a1a2e}\definecolor{dtflight}{HTML}{0f3460}
\definecolor{boxblue}{HTML}{e8f4f8}\definecolor{knownblue}{HTML}{1a3a6e}
\definecolor{cbKnownBg}{HTML}{E9E9FA}\definecolor{cbKnownFr}{HTML}{8E97C9}
\definecolor{cbNodeBg}{HTML}{F4F5F8}\definecolor{cbNodeFr}{HTML}{CDD2DB}
\definecolor{cbAssumeBg}{HTML}{FBEBEB}\definecolor{cbAssumeFr}{HTML}{C0392B}
\definecolor{cbResultBg}{HTML}{FBF4D2}\definecolor{cbResultFr}{HTML}{C8A21E}
\definecolor{cbTagBg}{HTML}{E2F1EF}\definecolor{cbTagFr}{HTML}{4FA295}
\definecolor{cbArrow}{HTML}{8A93A6}\definecolor{cbPyBg}{HTML}{ECEFF4}

\tcbset{
  keybox/.style={enhanced,breakable,colback=boxblue,colframe=dtfblue,boxrule=0.8pt,
    arc=3pt,left=10pt,right=10pt,top=8pt,bottom=8pt},
  cbflow/.style={enhanced,arc=3.5pt,boxrule=0.7pt,left=12pt,right=12pt,top=10pt,bottom=10pt,
    before skip=4pt,after skip=4pt},
}
\newcommand{\chainarrow}{\vspace{1pt}\centerline{\tikz{\draw[-{Latex[length=2.6mm,width=2.2mm]},
  line width=1pt,cbArrow](0,0.28)--(0,0);}}\vspace{1pt}}
\newcommand{\known}[2]{\begin{tcolorbox}[cbflow,colback=cbKnownBg,colframe=cbKnownFr]
  \centering\(\displaystyle #1\)\\[3pt]{\color{knownblue}\small[\,\textit{known: #2}\,]}\par
  \nobreak\vspace{1pt}\hfill{\tiny\color{gray}Known}\end{tcolorbox}}
\newcommand{\assume}[1]{\chainarrow\begin{tcolorbox}[cbflow,colback=cbAssumeBg,colframe=cbAssumeFr]
  {\color{cbAssumeFr}\bfseries[DTF input:]\ }{\color{cbAssumeFr}\small #1}\par
  \nobreak\vspace{1pt}\hfill{\tiny\color{gray}DTF input}\end{tcolorbox}}
\newcommand{\nodetext}[1]{\chainarrow\begin{tcolorbox}[cbflow,colback=cbNodeBg,colframe=cbNodeFr]
  \small #1\par\nobreak\vspace{1pt}\hfill{\tiny\color{gray}Derivation}\end{tcolorbox}}
\newcommand{\dtfresult}[2]{\chainarrow
  \begin{tcolorbox}[cbflow,colback=cbResultBg,colframe=cbResultFr,boxrule=1.3pt]
  \centering\(\displaystyle #1\)\par\nobreak\vspace{1pt}\hfill{\tiny\color{gray}DTF result}\end{tcolorbox}\chainarrow
  \begin{tcolorbox}[cbflow,colback=cbTagBg,colframe=cbTagFr]
  \centering{\itshape\color{dtflight}#2}\par\nobreak\vspace{1pt}\hfill{\tiny\color{gray}Takeaway}\end{tcolorbox}}
\newcommand{\pychip}[2]{{\footnotesize\colorbox{cbPyBg}{\ttfamily\,#1\,}\,{\color{gray}$\to$ #2}}}
\newcommand{\resultpy}[3]{\chainarrow
  \begin{tcolorbox}[cbflow,colback=cbResultBg,colframe=cbResultFr,boxrule=1.3pt]
  \centering\(\displaystyle #1\)\par\smallskip\centering\pychip{#2}{#3}\par
  \nobreak\vspace{1pt}\hfill{\tiny\color{gray}DTF result --- verified}\end{tcolorbox}}
\newcommand{\oneobject}[2]{\begin{tcolorbox}[enhanced,colback=dtfblue,colframe=dtfblue,arc=3pt,
  coltext=white,left=10pt,right=10pt,top=7pt,bottom=7pt,before skip=4pt,after skip=4pt]
  {\bfseries One object, seen here as: #1}\\[2pt]{\small #2}\end{tcolorbox}}
\newenvironment{chainbox}{\par\addvspace{4pt}\noindent\ignorespaces}{\par\addvspace{4pt}}
\newcommand{\chainhead}[2]{\subsection*{Chain #1\quad\normalfont\itshape #2}}
\newcommand{\setup}[1]{\noindent\emph{#1}\par\smallskip}
\newcommand{\lp}{\ell_P}

\hypersetup{colorlinks=true,linkcolor=knownblue,citecolor=knownblue,
  pdftitle={The Distorted Time Framework, Volume II: The Shape of Time},
  pdfauthor={William de Dufour}}
\pagestyle{fancy}\fancyhf{}
\fancyhead[L]{\small\itshape Distorted Time Framework --- Volume II}
\fancyhead[R]{\small\itshape The Shape of Time}
\fancyfoot[C]{\small\thepage}\renewcommand{\headrulewidth}{0.4pt}

\begin{document}

\begin{center}
{\large\bfseries The Distorted Time Framework\\[3pt]
\normalsize Volume II --- \itshape The Shape of Time}\\[6pt]
{\small Gravity, inertia, and the resolution of the black-hole singularity}\\[10pt]
{\normalsize William de Dufour}\\[2pt]
{\small\itshape Independent Researcher, Austin, TX \;\textrm{$\cdot$}\; \texttt{billd.dtf@gmail.com}}\\[3pt]
{\small July 2026 \quad---\quad Version 8}\\[3pt]
{\small\href{https://doi.org/10.5281/zenodo.21659853}{\texttt{DOI: 10.5281/zenodo.21659853}}}
\end{center}

\begin{abstract}
With time's local rate $u$ as the single primitive (Volume~I), gravity is the shape of that rate: a slow
body falls down its gradient, and the field law $\nabla^2u=4\pi G\rho/c^2$ is the rate's, not a graviton's.
The reconstructed metric matches general relativity in every tested regime (the Isenberg--Wilson--Mathews
conformally-flat scheme, exact to 1PN), and the frame field radiates exactly two tensor polarisations --- so
the tested content is GR's, and by Deser's uniqueness the two are indistinguishable wherever the spin-2
field is defined. The volume's distinctive claim lies where that field is \emph{not} defined: at Planck
curvature the emergent description of time \textbf{de-emerges}, replacing the singularity with a finite core
and a stable Planck-mass remnant --- a \emph{boundary conjecture} whose signatures (no echoes, no bounce)
differ sharply from the fuzzball and the Planck star. The rotating exterior is matched to Kerr (forward
through $O(a^3)$, all orders inherited via Deser). Stated plainly: the singularity results sit at the
epistemic level of a boundary conjecture, not a derived dynamics.
\end{abstract}

\smallskip
\noindent{\small\textbf{Keywords:} black-hole singularity; time de-emergence; general relativity; Kerr
metric; gravitational-wave polarisations; conformal flatness.}

\bigskip

\begin{tcolorbox}[keybox,unbreakable]
\textbf{The three commitments (from the Core, Volume~I).}\enspace This volume builds on three, argued there
and used here without re-argument (full recap: Appendix~A):
\begin{itemize}[leftmargin=1.4em,itemsep=1pt,topsep=3pt]
  \item \textbf{1.\ Time is primitive} --- its local rate $u(\mathbf x)$ is the one physical thing:
    Row~0 (timeless substrate, $\hat H\Psi=0$, energy conserved as a constraint) $\to$ Row~1 (the clock,
    where gravity and inertia live) $\to$ Row~2 (space).
  \item \textbf{2.\ Space is time's radiation} --- $\lambda_C=c\,T_C$; time takes $g_{00}=-u^2$ and space
    the complement $\gamma_{ij}=(2-u^2)\delta_{ij}$, while $u$ itself does \emph{not} propagate
    ($\pi_u\equiv0$). Hence zero new degrees of freedom, exact Lorentz ($\alpha_1=\alpha_2=0$), and a
    radiative content of exactly two --- read forward from what space is made of (Appendix~A). These carry
    the polarisation count and the singularity alike.
  \item \textbf{3.\ There is a proper now} --- the slice that solves Row-0's energy constraint
    (constant-$u$/CMC), real but empirically silent.
\end{itemize}
\noindent\emph{Here} the rate is \emph{sourced} by matter (gravity) and pushed to where it \emph{de-emerges}
(the singularity). What stays contingent is only the conformally-flat scheme this volume uses to compute the
metric beyond 1PN --- a limit on calculation, not on the commitments.
\end{tcolorbox}

\medskip
\noindent General relativity is our most successful theory of gravity, and it predicts its own failure:
at the centre of a black hole the curvature diverges and the theory stops. This volume says \emph{what
that failure is}. With time's rate as the primitive, the singularity is not a place where curvature
becomes infinite --- it is where the emergent description of time \textbf{de-emerges}, at Planck
curvature, and ``what is at $r=0$'' is a question asked in a language that has ceased to exist there.
What replaces the singularity is a finite core, and it makes predictions that differ, sharply, from
every other proposal. First, though, the gravity the singularity lives in.

\clearpage
\chainhead{1}{Gravity is the shape of the clock field}
\setup{A slow body falls down the gradient of time's rate. Make that rate primitive, and gravity is its shape --- with no force carrier.}
\begin{chainbox}
\known{g_{00}=-u^2,\qquad \ddot{\mathbf x}=-c^2\nabla u\ \ (\text{slow body})}
{Newtonian limit of GR: gravity is the curvature of \emph{time}}
\assume{promote $u$ to primitive; the one reparametrization-invariant action carries no $u$ time-derivative, so $\pi_u\equiv0$.}
\dtfresult{\nabla^2 u=\tfrac{4\pi G}{c^2}\rho\quad(\text{elliptic constraint, not a wave})}
{Gravity is the shape of the clock-rate field, sourced instantaneously --- no graviton to exchange. And
promoting time adds no degree of freedom and breaks no symmetry (Volume~I, Appendix~A): the radiation is
space, time does not radiate, $c$ is universal.}
\end{chainbox}

\chainhead{2}{General relativity, matched in the tested regime}
\setup{The reconstructed metric is GR's in every tested regime --- by construction, not by tuning; beyond it, matched only where the scheme is known to hold.}
\begin{chainbox}
\known{\text{reconstruct } g_{\mu\nu}\ \text{from } u\ \text{and the frame field}}{constrained, conformally-flat form}
\nodetext{this construction --- a conformally-flat spatial 3-metric plus elliptic constraints --- is not
a bespoke DTF scheme; it \emph{is} the Isenberg--Wilson--Mathews (IWM/CFC) approximation to
GR~\cite{Isenberg2008,WilsonMathews1995}, a known, well-posed, numerically validated formulation with a documented error
budget: exact for spherical symmetry, correct through 1PN (hence $\gamma=\beta=1$ falls out, not is tuned),
departs from GR at 2PN, and admits no conformally-flat slicing of Kerr at all. Naming the scheme converts an
open-ended ``closure item'' into a quantified one.}
\dtfresult{\gamma=\beta=1,\quad c_{\rm GW}=c\quad(\text{IWM/CFC, exact to 1PN})}
{The PPN parameters~\cite{Will2014} and the GW speed come out as GR's at the order IWM/CFC is known to
reproduce it. Honest caveat: the spatial metric is reconstructed in this \emph{conformally-flat
approximation} with a documented $g_{rr}$ deviation beyond 1PN (the closure item, sharpened above) --- so
this is GR \emph{proposed as what the clock field is made of} at the tested level, not a claim of exact
recovery. The binary-pulsar orbital decay --- matched to $\sim0.2\%$ against
PSR~B1913+16~\cite{WeisbergHuang2016} --- is a \emph{radiation-reaction} effect, which the \emph{waveless}
conformally-flat scheme does not itself produce: it is GR's quadrupole result, and DTF inherits it not
through the CFC reconstruction but through the \emph{Deser closure} (box after Chain~8) --- the radiative
sector is a two-polarisation spin-2 field, so its matter coupling, and hence the $G/5c^5$ quadrupole
emission, are GR's (a mode \emph{count} alone would not fix the prefactor). A genuinely independent DTF
quadrupole computation from the frame field's own dynamics is not written out in this trilogy.}
\end{chainbox}

\clearpage
\chainhead{3}{Space radiates --- two tensor polarisations}
\setup{The radiative sector is space. How many ways it can wave follows from what space is made of --- not from a projector chosen in advance.}
\begin{chainbox}
\known{h_{ij}\ \text{symmetric}:\ \tfrac{d(d+1)}{2}=6\ \text{components at}\ d=3}{the extent structure --- what assigns a length to each direction}
\nodetext{space is compiled from extents (Volume~I, Chain~5), so the object that can modulate is $h_{ij}$, and the
frame's three internal rotations are redundancy in it (they leave $g_{ij}=e^a_ie^a_j$ fixed at all orders).
Split the six by propagation direction --- $2$ scalar, $2$ vector, $2$ tensor --- and the framework removes
four of them: the scalars because a uniform rescaling of extents \emph{is} a clock-rate change, which the
radiation rule forbids from propagating (Volume~I, Chain~5) and the gravity field law fixes elliptically
(Chain~1); the vectors because space is compiled, so relabelling points within a slice is not a physical change.}
\resultpy{6-2\,(\text{scalar})-2\,(\text{vector})=2\ \text{transverse-traceless}}
{DTF\_radiative\_content\_constructive.py}{2, forward from the premises}
\dtfresult{\text{exactly two tensor polarisations, no breathing mode}}
{This is the framework's cleanest falsifiable stake: a detected scalar/breathing gravitational-wave
polarisation would refute it. Time (the scalar $u$) does not radiate --- and the count is now read forward
from what space is, with the familiar TT projector appearing as a \emph{consequence} of the split rather
than as its premise (Volume~I, \S A.4).}
\end{chainbox}
\oneobject{inertia \& gravity}{One rate $\omega_C$: it \emph{is} the inertia (linear in $m$,
definitional), and it \emph{sources} the gravitational well in $u$ (the $G$-coupled $\sim(m/m_P)^2$
dent), whose extremum is free fall. The equivalence principle is that single identity --- read forward
into the chain below, where inertia is the rate, not a wake.}

\clearpage
\chainhead{4}{Inertia --- the rate, not a wake}
\setup{Inertia needs no mechanism. Mass \emph{is} the clock's ticking rate; inertia is that rate, carried by a moving body.}
\begin{chainbox}
\known{\mathbf p=\nabla S=\gamma m\mathbf v,\qquad E=-\partial_t S=\gamma mc^2}{Hamilton--Jacobi: the moving clock's momentum and energy}
\dtfresult{\text{inertia}=\text{the rate }\omega_C\ \text{itself}\ (m=\hbar\omega_C/c^2),\ \text{linear in }m,\ \text{no }G}
{There is no self-energy to compute. Mass \emph{is} the ticking rate (Volume~I, ``a note on mass''); inertia
is that rate carried by a moving body --- linear in $m$, $\hbar$-set, carrying no $G$. The coefficient
$mc^2$ is not a debt but a unit conversion ($E=\hbar\omega_C$). A self-field reaction to the body's own
$u$-well could not be the source: that well is $G$-coupled and $\sim(m/m_P)^2$ shallow --- some
forty-five orders too weak to be inertia, which in any case needs no such mechanism. What remains owed is
only the \emph{pattern} of ticking rates (the mass hierarchy) --- a free parameter for every theory alike;
deferred, not claimed.}
\end{chainbox}

\clearpage
\chainhead{5}{The singularity: time de-emerges}
\setup{Curvature does not diverge --- the description does. Where the clock can no longer complete a tick, Row-1 time ends.}
\begin{chainbox}
\known{K=R_{\mu\nu\rho\sigma}R^{\mu\nu\rho\sigma}=\frac{48\,(GM/c^2)^2}{r^6},\quad L_{\rm curv}=K^{-1/4}}
{Kretschmann invariant, Schwarzschild}
\assume{Row-1 time de-emerges where the clock loses definiteness: one coherent tick no longer fits in a
curvature time, $N\equiv L_{\rm curv}/\lp\to1$ (Planck curvature, coefficient-free). The Planck-curvature
\emph{threshold} is imported --- most quantum-gravity programmes agree curvature saturates there; what
is \emph{native} is the interpretation: de-emergence of Row-1 time, not a bounce or a new phase.}
\resultpy{r_*=48^{1/6}\lp\,(M/m_P)^{1/3},\quad \frac{r_*}{r_s}\simeq0.95\,(M/m_P)^{-2/3}}
{DTF\_singularity\_deemergence.py}{deep interior, fat core}
\dtfresult{\text{a finite core at } r_*,\ \text{never the point } r=0}
{The de-emergence surface sits parametrically \emph{inside} the horizon ($\sim10^{-26}\,r_s$ for a solar
mass); the horizon stays curvature-regular. The singularity is a question in a language that has ceased
to exist --- and what is left at the core is Row~0 showing through: the timeless substrate, uncovered
where Row-1 time ends. The ontology closes on itself.}
\end{chainbox}

\clearpage
\chainhead{6}{The interior: a clock that halts before the crunch}
\setup{Read through a geodesic comoving clock, collapse stops a Planck time short of the classical singularity.}
\begin{chainbox}
\known{\rho(\tau_{\rm rem})=\frac{1}{6\pi G\,\tau_{\rm rem}^2}}{matter-dominated collapse, comoving proper time}
\assume{the same criterion in density: de-emergence at $\rho\to\rho_P=c^5/\hbar G^2$ (Planck density).}
\resultpy{\tau_{\rm rem}^*=\frac{t_P}{\sqrt{6\pi}}\approx0.23\,t_P\ \text{before the crunch}}
{DTF\_LTB\_collapse.py}{mass- \& profile-indep.}
\dtfresult{\text{halt at finite scale factor, never } a=0}
{The comoving clock reaches $\delta u/u\sim1$ a universal $0.23\,t_P$ before the classical singularity ---
independent of mass and (for inhomogeneous LTB dust~\cite{LTB}) of the density profile. A fat Planck-density core,
not a point.}
\end{chainbox}

\clearpage
\chainhead{7}{Observables --- how the finite core differs from its rivals}
\setup{A finite core in place of a singularity has consequences at both ends of a black hole's life.}
\begin{chainbox}
\known{r_s\propto M,\qquad R_*=(M/\rho_P)^{1/3}\propto M^{1/3}}{horizon and core shrink at different rates}
\resultpy{M_{\rm rem}=\frac{m_P}{\sqrt8}\approx0.35\,m_P,\quad \text{no echoes},\ \text{no bounce}}
{DTF\_singularity\_observable.py}{stable Planck-mass relic}
\dtfresult{\text{a stable, horizonless Planck-mass remnant}}
{Evaporation halts when the core meets the shrinking horizon --- no complete evaporation, no final burst.
The structure sits at $\sim10^{-26}r_s$, so \emph{no} post-ringdown echoes; time de-emerges rather than
bouncing, so \emph{no} radio/gamma bounce transient. Whether a population of such relics could contribute to
dark matter is a separate question this volume does not assess.}
\end{chainbox}

\begin{center}\small
\renewcommand{\arraystretch}{1.2}
\begin{tabular}{@{}L{5.0cm}L{3.6cm}L{4.2cm}@{}}
\toprule
\textbf{Signature} & \textbf{DTF (de-emergence)} & \textbf{Rivals}\\
\midrule
Post-ringdown GW echoes & none (regular horizon) & fuzzball/firewall\cite{Mathur2005,AMPS2013}: yes\\
Evaporation endpoint & Planck remnant $\sim0.35\,m_P$ & complete / burst\\
Primordial-hole fate & stable Planck-mass relic & Planck star\cite{RovelliVidotto2014}: bounces\\
Radio/gamma bounce transient & none & Planck star: yes\\
\bottomrule
\end{tabular}
\end{center}

\clearpage
\chainhead{8}{The rotating sector --- Kerr matched (exterior; forward through $O(a^3)$, higher orders inherited)}
\setup{Rotation is the one place the conformally-flat reconstruction stops. Sourced forward from DTF's own sectors --- not read off Kerr --- it is reached even so: frame-dragging is the proper now dragged, and the spatial geometry follows to second order in spin.}
\begin{chainbox}
\known{\mathrm{tr}\,K=-\tfrac{1}{N}D_i\beta^i=0}{a stationary axisymmetric const-$u$ slice is \emph{maximal} --- a legitimate proper now}
\assume{rotation enters only through the shift $\beta^i$, sourced by the matter momentum $\rho v^i$ \emph{alone} ($\pi_u\equiv0$: no clock current to be dragged alongside it).}
\resultpy{\omega=\frac{2GJ}{c^2r^3}\ (\text{DTF, forward})=\omega_{\rm Kerr}^{\rm far},\qquad N_{\rm dof}=2}
{DTF\_kerr\_constructive.py}{frame-dragging is native; rotation adds no mode}
\end{chainbox}
\begin{chainbox}
\known{C_{ijk}=0\ \Longleftrightarrow\ \text{conformally flat}\ \Longleftrightarrow\ \text{no physical tensor content}}
{the coordinate-invariant obstruction (Cotton tensor)}
\nodetext{compute $C_{ijk}$ of the Kerr spatial 3-geometry on that slice, order by order in spin: it is
zero at $O(a^0)$ (Schwarzschild) and $O(a^1)$ (frame-dragging does not curve space), and \emph{first
nonzero at $O(a^2)$}. So the tensor sector must switch on exactly there --- and its source is the
\emph{quadratic} of the $O(a)$ shift, a gravitomagnetic dipole (hence angular content $\ell=0,2$).}
\resultpy{\delta R[g_2]=-\,\delta^2R[g_1]\quad\text{balances to}\ 10^{-126}}
{DTF\_kerr\_oa2\_sufficiency.py}{frame-drag$^2$ sources the $O(a^2)$ TT mode}
\dtfresult{\textbf{DTF}=\textbf{Kerr through }O(a^2),\ \text{coefficient-exact}}
{The frame-drag$^2$ source, pushed through DTF's radiative tensor operator, reproduces Kerr's
second-order geometry \emph{exactly}; by uniqueness of the elliptic problem the forward solve \emph{is}
Kerr. This is proven explicitly here at $O(a^2)$; the extension to \emph{all orders} is \emph{inherited}
from the Deser bootstrap (box below), not computed order by order. It is
the \emph{exterior} rotating geometry (frame-dragging, ergosphere $|\beta|=uc$); rotating collapse is a
separate question, not claimed here.}
\end{chainbox}

\begin{tcolorbox}[keybox,unbreakable]
\textbf{What the rotating sector rests on --- and why it closes to all orders.}
\begin{itemize}[leftmargin=1.5em,itemsep=3pt,topsep=3pt]
  \item \textbf{The reframe.} Conformal flatness is the \emph{computational} (IWM/CFC) scheme, not the
    ontology. Kerr admits a maximal ($K=0$) slice --- the very slice DTF calls the proper now ---
    and Garat--Price~\cite{GaratPrice2000} forbids only a \emph{conformally-flat} one. The theorem bounds
    the old reconstruction, not the framework.
  \item \textbf{The operator is not a premise --- it is Lorentz invariance.} That DTF's \emph{stationary}
    tensor operator equals its \emph{radiative} one is not an extra assumption: universal $c$ (Volume~I,
    Chain~5) makes the wave operator $\Box=-c^{-2}\partial_t^2+\nabla^2$ the \emph{only} second-order
    operator the transverse-traceless frame field can obey. The linearised field operator on a TT mode is
    exactly $-\tfrac12\Box h_{ij}$ (two propagating polarisations), and its stationary reduction
    ($\partial_t\!\to\!0$) is $-\tfrac12\nabla^2$ --- the very operator that balanced the frame-drag$^2$
    source. Radiative and near-zone are one operator, evaluated at two values of $\partial_t$.
    \pychip{DTF\_kerr\_premise\_oneoperator.py}{$\delta R[h^{\rm TT}]=-\tfrac12\Box h$, stationary $=\nabla^2$}
  \item \textbf{Native through $O(a^3)$; all orders by Deser --- and the two are graded.} The operator
    ($\Box$, all orders) and the quadratic vertex (the radiative milestone) fix the solution \emph{through}
    $O(a^2)$; $O(a^3)$ is the first order whose source is \emph{cubic} (frame-drag sector, by parity), probing
    the cubic vertex. That vertex is not free: DTF's frame field is established to be a two-polarisation
    spin-2 metric (App.~A / Volume~I) whose internal $SO(3)$ is gauge (invisible in $g_{ij}=e^a_ie^a_j$ to all
    orders) and whose clock $u$ is non-propagating ($\pi_u\equiv0$), so its physical propagating content is
    precisely a consistent, Lorentz-invariant, two-derivative, two-polarisation spin-2 field --- for which
    Deser's self-coupling bootstrap~\cite{Deser1970} makes general relativity the \emph{unique} completion,
    every vertex GR's. \textbf{The grading matters, and we state it:} the $O(a^2)$ and $O(a^3)$ agreements are
    \emph{computed forward} in DTF (closing to $10^{-126}$ and $10^{-131}$); the extension to \emph{all}
    orders is \emph{inherited} from Deser's theorem --- which any two-polarisation spin-2 theory shares ---
    not an order-by-order DTF calculation. Granted that inheritance, DTF $=$ Kerr to all orders in the
    exterior, and the rotating sector adds no distinctive signature (the distinctive stakes remain the GW
    polarisations and the black-hole cores). The one honest residual is the rotating \emph{interior}
    (rotating collapse / de-emergence), untouched here.
    \pychip{DTF\_kerr\_deser\_closure.py}{native $\le O(a^3)$; all orders inherited via Deser (2-pol spin-2)}
\end{itemize}
\end{tcolorbox}

\begin{tcolorbox}[keybox]
\textbf{What Deser closes, and what it does not.}\enspace The bootstrap~\cite{Deser1970} is a theorem about
a \emph{field}: any consistent, Lorentz-invariant, two-derivative, two-polarisation spin-2 field completes
uniquely to general relativity. DTF's frame field is such a field, so wherever that field is defined DTF's
vertices are GR's --- which is what closes the rotating exterior to all orders, and \emph{equally} what
makes the gravitational sector observationally indistinguishable from GR in every tested regime. It follows
that no result of this volume's core sector can be a modification of GR's field equations, and none is
offered as one. De-emergence is a claim about \emph{where the spin-2 description ceases to be defined} ---
at $\delta u/u\sim1$ --- and Deser's theorem, like every perturbative uniqueness result, is silent there.
This places the finite core, the halt, and the remnant at the epistemic level of a \emph{boundary
conjecture}: the same level as the Planck star and the fuzzball, differing in \emph{what} is claimed to lie
at the boundary, not in derivational status. The comparison table (Chain~7) is a comparison of boundary
conjectures by signature, not a derivation discriminating against rivals.
\end{tcolorbox}

\begin{tcolorbox}[keybox,unbreakable]
\textbf{What it stakes, and what it owes.}
\begin{itemize}[leftmargin=1.5em,itemsep=2pt,topsep=3pt]
  \item \textbf{Falsifiable stakes, and which discriminate.} Two tensor GW polarisations, no
    breathing mode --- a falsifiable stake, but one \emph{shared with} GR, not a discriminator. The
    echo-free finite core, the stable Planck-mass remnant, and the absence of a bounce distinguish DTF's
    boundary conjecture from rival \emph{conjectures} (fuzzball, Planck star) \emph{by signature} --- a
    comparison of conjectures at one epistemic level, not a derivation ranking DTF above them (see the Deser
    box after Chain~8).
  \item \textbf{Honest limits.} What is removed is the \emph{curvature} singularity and the classical
    endpoint; whether the object is geodesically \emph{complete} --- as opposed to a worldline exiting
    Row~1 --- is a separate, open question. Every numerical coefficient ($r_*$, the $0.23\,t_P$ halt,
    $M_{\rm rem}$) is an $O(1)$ illustration in a tractable model; the load-bearing claim (de-emergence at
    Planck curvature) is coordinate-invariant and coefficient-free.
  \item \textbf{Rotation: reached (Chain~8), not excluded.} That Kerr admits no conformally-flat
    slicing~\cite{GaratPrice2000} bounds only the IWM/CFC \emph{reconstruction}, not the ontology:
    Kerr does admit a maximal ($K=0$) slice --- DTF's
    proper now --- and sourced forward on it, DTF reproduces Kerr's frame-dragging natively (with $N_{\rm
    dof}=2$) and its full $O(a^2)$ geometry coefficient-exactly (Chain~8), the tensor operator being fixed by
    Lorentz invariance rather than assumed. And it closes: because the physical frame field is a
    two-polarisation spin-2 metric ($SO(3)$ internal gauge, $\pi_u\equiv0$), Deser's uniqueness forces every
    vertex to GR's, so DTF $=$ Kerr to \emph{all orders} in the exterior (computed forward at $O(a^2)$ and
    $O(a^3)$; the rest inherited from Deser, not computed order-by-order).
    What remains open is only the \emph{interior} of a rotating hole (rotating collapse / de-emergence),
    which Chain~8 does not claim.
  \item \textbf{Owed.} The $O(1)$ core coefficient; an anisotropic (BKL) illustration; the full nonlinear
    closure of $S_u+S_m+S_{\rm frame}$. (\emph{Not} owed: the magnitude of inertia ---
    mass \emph{is} the ticking rate, so there is no self-energy to compute; only the dimensionless rate
    \emph{ratios}, the hierarchy, remain, and those every theory owes alike. Chain~4.)
\end{itemize}
\noindent \textbf{Volume~I} builds the ontology; \textbf{Volume~III} reads the quantum. The three volumes
are one \emph{organizing principle} (Volume~I) checked in three sectors --- and this volume's finite core
is not a separate result but the \emph{same} clock-definiteness surface as quantum measurement
(Volume~III), crossed in the opposite direction: there a clock \emph{becomes} definite, here it
\emph{de-emerges}. The singularity and the measurement event are one transition, run both ways. A fourth
paper --- deep cosmology (the proper now as cosmic time) --- is in preparation.
\end{tcolorbox}

\begin{tcolorbox}[keybox,unbreakable]
\textbf{Provenance of the checks (ontology-first).}\enspace This volume carries the trilogy's genuinely
\textbf{DTF-native} computational content: the de-emergence surface and fat-core scaling
(\texttt{DTF\_singularity\_deemergence}), the profile-independent $0.23\,t_P$ collapse halt
(\texttt{DTF\_LTB\_collapse}), and the remnant / no-echo observables
(\texttt{DTF\_singularity\_observable}) --- real results, with Hawking evaporation~\cite{Hawking1975}
invoked only as an external Row-2 input. The rotating sector (Chain~8) is likewise \textbf{computed forward}
--- though \emph{confirmatory rather than discriminating}, since the Deser closure already fixes the
vertices, so these agreements could not have come out otherwise; their value is as a self-test of the
construction, and it passes (to $10^{-126}$/$10^{-131}$): the frame-dragging coefficient (\texttt{DTF\_kerr\_constructive}), the order-by-order Cotton
census on the maximal slice (\texttt{DTF\_kerr\_forward\_tensor}, self-tested to give zero at $O(a^0)$ and
$O(a^1)$), the coefficient-exact $O(a^2)$ and $O(a^3)$ balances $\delta R[g_n]=-\,\delta^{\ge2}R[g_{<n}]$
(\texttt{DTF\_kerr\_oa2\_sufficiency}, \texttt{DTF\_kerr\_oa3\_balance}), the one-operator identity
$\delta R[h^{\rm TT}]=-\tfrac12\Box h$ (\texttt{DTF\_kerr\_premise\_oneoperator}), and the Deser closure
(\texttt{DTF\_kerr\_deser\_closure}) --- Kerr entering only as the target compared against, never as an
input to the derivation. The provenance is graded honestly: the rotating sector is computed
forward \emph{through $O(a^3)$} (confirmatory, per above); the extension to all orders is \emph{inherited} from the Deser bootstrap
(shared by any two-polarisation spin-2 theory), a theorem, not an order-by-order DTF computation. The
gravity chains (1--4) are \textbf{consistency}: GR is \emph{proposed as} what
the clock field is made of (conformally-flat-approximate, $g_{rr}$ open), and the two-polarisation count
is standard linearised gravity. Nothing here rests on circular reasoning.
\end{tcolorbox}

\appendix
\section*{Appendix A --- The DTF Core in brief}
\addcontentsline{toc}{section}{Appendix A}
\noindent For readers coming to this volume first. The Distorted Time Framework takes the local
\emph{rate} of time as the one physical primitive, and builds outward in three layers (Volume~I).
\begin{itemize}[leftmargin=1.4em,itemsep=3pt,topsep=3pt]
  \item \textbf{Row 0 --- the timeless substrate.} $\hat H\Psi=0$, no external time; correlations
    simultaneous, not ordered.
  \item \textbf{Row 1 --- the clock.} A definite clock-rate field $u(\mathbf x)>0$; its reading
    $\tau=\int u\,dt$ accumulates monotonically --- time, and its arrow. Gravity, inertia, and phase live
    here.
  \item \textbf{Row 2 --- space.} An extent is a period radiated at $c$ ($\lambda_C=c\,T_C$): space is
    time's radiation, compiled by a frame field. Because $c$ is one universal exchange rate, Lorentz
    invariance is its signature.
  \item \textbf{One object, four faces.} Extremising the interval (Fermat, applied to time) gives one
    phase read four ways: inertia ($\mathbf p=\nabla S$), free fall (geodesic), the pilot wave
    ($\mathbf v=\nabla S/m$), and the path integral. Inertia (the rate $\omega_C$) and the pilot wave (the
    space-face gradient $\nabla S$) are tied by $\mathbf p=\nabla S=m\mathbf v$, both $\hbar$-set; the
    $G$-coupled well the rate digs is gravity, a separate object.
  \item \textbf{Non-propagating time.} $u$ enters the action with no time-derivative, so $\pi_u\equiv0$:
    zero propagating degrees of freedom, and a preferred slice (the \emph{proper now}) that is real but
    empirically inert. Space radiates (two tensor polarisations, read forward from the extent structure); time does not (no breathing mode).
  \item \textbf{Definiteness.} A single clock-definiteness threshold governs the black-hole core (this
    volume) and quantum measurement (Volume~III).
\end{itemize}

\section*{Code and data availability}
\addcontentsline{toc}{section}{Code and data availability}
{\small The SymPy/NumPy verification scripts named throughout (\texttt{DTF\_singularity\_deemergence},
\texttt{DTF\_LTB\_collapse}, \texttt{DTF\_singularity\_observable}, and the six-script Kerr suite named in
the Provenance box --- \texttt{DTF\_kerr\_constructive}, \texttt{DTF\_kerr\_forward\_tensor},
\texttt{DTF\_kerr\_oa2\_sufficiency}, \texttt{DTF\_kerr\_oa3\_balance},
\texttt{DTF\_kerr\_premise\_oneoperator}, \texttt{DTF\_kerr\_deser\_closure}) are publicly available at
\url{https://github.com/billd-dtf/Distorted-Time-Framework} (release \texttt{v1.0}, matching this
revision); see the repository's README for a script-to-claim index and run instructions.\par}

\section*{Author's Note: the role of AI}
\addcontentsline{toc}{section}{Author's Note: the role of AI}
{\small
The Distorted Time Framework---its central hypothesis (time's local \emph{rate}, the ADM lapse,
as the one physical primitive), its ontology-first strategy, and the physical
interpretations---was conceived and developed by the author over several years of independent
research outside academic institutions. The mathematical formalization was carried out in
collaboration with Claude (Anthropic): the gravitational field equations and degree-of-freedom
count, the black-hole singularity de-emergence and Planck-remnant calculations, the rotating-sector
(Kerr) derivations, and the symbolic/numerical verification scripts (SymPy/NumPy) referenced
throughout. Separately, other large language models (ChatGPT, Grok, Gemini) were used as independent
referees, providing adversarial review of the manuscript; their critiques shaped the revisions but
not the derivations.

In each case the author specified the physical target, the constraints the result had to satisfy,
and the direction of the argument, and reviewed and approved each result; the AI executed and
checked the formal derivations. Without that assistance the author could not have produced the
mathematical formalization at this level of detail.

This is disclosed because intellectual honesty requires it. It does not change the epistemic status
of the results---a derivation is correct or it is not, independent of how it was produced---but it
does mean they warrant careful independent scrutiny rather than assumed authority. The author
welcomes that scrutiny; the explicit claim-type labels and the reproducible verification scripts are
provided precisely to enable it.
\par}

\begingroup\small
\begin{thebibliography}{9}
\addcontentsline{toc}{section}{References}
\bibitem{Hawking1975} S.~W.~Hawking, ``Particle Creation by Black Holes,''
  \emph{Commun.\ Math.\ Phys.} \textbf{43}, 199--220 (1975).
\bibitem{WeisbergHuang2016} J.~M.~Weisberg and Y.~Huang, ``Relativistic Measurements from Timing the
  Binary Pulsar PSR~B1913+16,'' \emph{Astrophys.\ J.} \textbf{829}, 55 (2016).
\bibitem{Will2014} C.~M.~Will, ``The Confrontation between General Relativity and Experiment,''
  \emph{Living Rev.\ Relativity} \textbf{17}, 4 (2014).
\bibitem{LTB} G.~Lema\^itre, \emph{Ann.\ Soc.\ Sci.\ Bruxelles} \textbf{A53}, 51 (1933);
  R.~C.~Tolman, \emph{Proc.\ Natl.\ Acad.\ Sci.\ USA} \textbf{20}, 169 (1934);
  H.~Bondi, \emph{Mon.\ Not.\ R.\ Astron.\ Soc.} \textbf{107}, 410 (1947).
\bibitem{Mathur2005} S.~D.~Mathur, ``The fuzzball proposal for black holes: an elementary review,''
  \emph{Fortsch.\ Phys.} \textbf{53}, 793--827 (2005).
\bibitem{AMPS2013} A.~Almheiri, D.~Marolf, J.~Polchinski, and J.~Sully, ``Black Holes:
  Complementarity or Firewalls?'' \emph{J.\ High Energy Phys.} \textbf{02}, 062 (2013).
\bibitem{RovelliVidotto2014} C.~Rovelli and F.~Vidotto, ``Planck stars,''
  \emph{Int.\ J.\ Mod.\ Phys.\ D} \textbf{23}, 1442026 (2014).
\bibitem{Isenberg2008} J.~Isenberg, ``Waveless approximation theories of gravity,''
  \emph{Int.\ J.\ Mod.\ Phys.\ D} \textbf{17}, 265--273 (2008) [1978 unpublished manuscript].
\bibitem{WilsonMathews1995} J.~R.~Wilson and G.~J.~Mathews, ``Instabilities in Close Neutron Star
  Binaries,'' \emph{Phys.\ Rev.\ Lett.} \textbf{75}, 4161--4164 (1995).
\bibitem{Deser1970} S.~Deser, ``Self-interaction and gauge invariance,''
  \emph{Gen.\ Relativ.\ Gravit.} \textbf{1}, 9--18 (1970).
\bibitem{GaratPrice2000} A.~Garat and R.~H.~Price, ``Nonexistence of conformally flat slices of the
  Kerr spacetime,'' \emph{Phys.\ Rev.\ D} \textbf{61}, 124011 (2000).
\end{thebibliography}
\endgroup

\end{document}
